\subsection{Morphisms of quivers}

Let C and C' be quivers. A morphism of quivers from C to C' is a pair \((u,v)\) , where \(u\colon\operatorname{Som}(\mathrm{C})\to\operatorname{Som}(\mathrm{C}')\) and \(v\colon\operatorname{Fl}(\mathrm{C})\to\operatorname{Fl}(\mathrm{C}')\) are maps such that \(o_{C'}\circ v=u\circ o_{C}\) and \(t_{C'}\circ v=u\circ t_{C}\) .

Let \(\varphi = (u, v)\) be a morphism of quivers from C to C'. The maps \(u\) and \(v\) are denoted \(\operatorname{Som}(\varphi)\) and \(\operatorname{Fl}(\varphi)\) . If \(a\) is a vertex of C, one will denote, by abuse of notation, \(\varphi(a)\) the image of the vertex \(a\) under \(\varphi\) and we will say that \(\varphi(a)\)

is the image of \(a\) under \(\varphi\) . Similarly, if \(f\) is an arrow of C, we will denote, by abuse of notation, \(\varphi(f)\) the arrow \(v(f)\) of C' and we will say that it is the image of \(f\) under \(\varphi\) .

Let C, \(C'\) , \(C''\) be three quivers, namely \(\varphi = (u, v)\) a morphism from C to \(C'\) and let \(\varphi' = (u', v')\) be a morphism from \(C'\) to \(C''\) . Then \((u' \circ u, v' \circ v)\) is a morphism from C to \(C''\) , denoted \(\varphi' \circ \varphi\) and which is called the composite of \(\varphi'\) and of \(\varphi\) .

We denote \(\mathrm{Id}_{\mathrm{C}}\) the morphism \((\mathrm{Id}_{\mathrm{Som}(\mathrm{C})}, \mathrm{Id}_{\mathrm{Fl}(\mathrm{C})})\) from C into itself.

Let \(\varphi\) be a quiver morphism from C to C'. For \(\varphi\) to be an isomorphism, it is necessary and sufficient that the maps \(\operatorname{Som}(\varphi)\) and \(\operatorname{Fl}(\varphi)\) are bijective (cf. E, IV, p. 6).

Thus, if one calls a quiver structure on sets S and F the data of two maps \(o, t: \mathrm{F} \to \mathrm{S}\) one can take quiver morphisms as morphisms of this structure (E, IV, p. 11).

Let \(\varphi\) be a quiver morphism from C to C'. There exists a unique subquiver of C' whose sets of vertices and arrows are the images of \(\mathrm{Som}(\varphi)\) and \(\mathrm{Fl}(\varphi)\) respectively. It is denoted \(\varphi(\mathrm{C})\) and it is called the image of C under \(\varphi\) . One says that \(\varphi\) is injective if the maps \(\mathrm{Som}(\varphi)\) and \(\mathrm{Fl}(\varphi)\) are injective; in this case, \(\varphi\) induces an isomorphism from C onto its image.

Let \(\varphi\) and \(\psi\) be morphisms of quivers from C to C'. There exists a unique subquiver of C whose vertices and arrows are respectively the vertices and arrows of C having the same image under \(\varphi\) and \(\psi\) . We call it the equalizer of \(\varphi\) and \(\psi\) .

